\section{Introduction}\label{sec:intro}

There are two ways to hold a share of a company for the next month. One is to buy the share today. The other is to buy a
futures contract that delivers the share at the end of the month. Since both end in the same place, their prices should
be linked: the futures price should equal today's share price, plus the interest one would pay to finance the share,
minus any dividend paid before delivery. If the future is dearer than this \emph{cost-of-carry} value, a trader can buy
the share, sell the future and hold both until delivery, earning the difference without any exposure to the share
price. This is \emph{cash-and-carry arbitrage}. If the future is cheaper, the reverse trade (sell the share, buy the
future) earns the difference. The gap between the futures price and its cost-of-carry value is called the
\emph{basis}. In a market where these trades are free and unlimited, the basis should be close to zero at all times.

In real markets arbitrage has limits. Every trade pays taxes and fees. The reverse trade needs a short sale, which may be
costly or not allowed. An arbitrage position ties up capital, and the gap can widen before it closes, so the
arbitrageur bears risk and may be forced to close early (Shleifer and Vishny, 1997; Gromb and Vayanos, 2010).
Exchanges also cap the positions traders may hold. With such limits, the basis can wander inside a band without anyone
finding it worth closing, and arbitrage pushes back only when the basis leaves the band. How wide this band is, and
which limits set its edges, are empirical questions.

India is an unusually good place to answer them, because most of its limits to arbitrage have changed on known dates.
The securities transaction tax (STT) on the sale of futures was raised in April 2023, October 2024 and April 2026. Only
stocks with derivatives, the stocks of the futures and options (F\&O) segment, may be sold short, and in practice the
overnight short position in an Indian stock is a short future. Arbitrage mutual funds, which hold exactly the
cash-and-carry trade, grew about fourfold after a change in the taxation of debt funds in 2023. And position limits are
set by formula and were redesigned in October 2025.

India also has a rule that, to my knowledge, has not been studied. Each stock's derivatives have a
\emph{market-wide position limit} (MWPL). When open interest, the number of contracts outstanding across all traders,
exceeds 95\% of this limit at the end of a day, the stock enters an \emph{F\&O ban} the next day. During the ban every
trader, not only arbitrageurs, may only trade to reduce positions, and the ban lasts until open interest falls to 80\%
of the limit. Between 2015 and October 2026 this happened 1,788 times, to 141 different stocks. A ban suddenly removes the
possibility of opening new arbitrage positions in a stock, which makes it a sharp test of what position limits do.

The project asks how strongly these limits bind in Indian equities and how much mispricing they explain. It has two
stages. This report covers the first. It sets out a model of arbitrage under India's limits, fixes the tests in advance,
builds the data, and answers three questions: how wide is the band within which the basis of Indian stock futures
moves; do trading costs and other frictions set its edges; and does the F\&O ban actually restrict trading in the way the
rule intends? The second stage asks what the ban and the other limits do to prices.

I proceed in a fixed order, designed to keep theory and data from being fitted to each other, which was the main advice
of my supervisor at the start of the project. I first wrote a model in which each limit enters as a parameter. From the
model I derived predictions, each with the result that would contradict it. I then wrote these tests into a
pre-analysis plan and froze it with a time stamp before I related any outcome to any treatment. The data, built from
public files of the National Stock Exchange (NSE), cover 486,040 stock-days of the futures-cash basis for 371 stocks.

The findings of this stage are the following.

\begin{enumerate}[leftmargin=1.4em,itemsep=3pt]
\item \textbf{The basis moves in a wide, asymmetric band.} A three-regime threshold model places the edges of the band
about 35 basis points (bp) below and 24\bp\ above each stock's typical basis. Inside the band the basis drifts; outside
it, it moves back. The lower edge is further away, as expected when the reverse trade, which needs a short sale, is
costlier. A test of the linear model against the threshold model rejects linearity decisively. Part of the day-to-day
reversion that the model measures is, however, noise in daily prices rather than arbitrage.
\item \textbf{Costs shape the band, but loosely.} The full cost of opening a new cash-and-carry position and holding it
to delivery is 27 to 31\bp, which places it in the outer tail of the basis, while rolling an existing position costs
only 4 to 8\bp. The upper part of the basis distribution rose after the 2024 and 2026 tax increases, as the model
predicts, and the test fixed in advance supports this. But the rise was larger than the change in costs and was not
unusual compared with the same statistic on other dates, so I read the evidence as weak. Across stocks, the band is
wider where trading frictions are higher.
\item \textbf{The ban binds hard on trading.} On ban days futures trading falls by about 87\% relative to the same
stock's days at similar levels of open interest, and cash-market trading moves toward delivery. How the rule was
enforced matters: once violations were judged on intraday rather than end-of-day positions, from March 2020, every
futures trade during a ban closed an existing position. There is no sign that traders hold back to avoid the trigger.
\item \textbf{The 2025 rules cut the number of bans.} Measured against the old rule applied to the same stocks, the
October 2025 rules produced about a quarter of the ban days in their first year.
\end{enumerate}

The closest earlier studies of Indian stock futures estimate no-arbitrage bands from intraday data over a few months
or years (Aggarwal, 2015; Shankar, Sankar and Kiran Kumar, 2019). This report adds a model of arbitrage that includes
India's position limits and the ban, estimates the band over twelve years of daily data for every stock with futures,
relates it to dated changes in costs, and gives the first evidence on how the F\&O ban affects trading.

\Cref{sec:setting} describes the Indian rules. \Cref{sec:lit} reviews the related literature. \Cref{sec:model} presents
the model and its predictions. \Cref{sec:data} describes the research design and the data. \Cref{sec:band} estimates the
band, \cref{sec:costs} asks what sets it, and \cref{sec:ban} studies trading in a ban. \Cref{sec:conclusion} concludes
and sets out the second stage. The appendices give the proofs and the details of the test plan and the data.
